\clearpage
\section*{Problem 6}
\addcontentsline{toc}{section}{Problem 6}
\begingroup
\hypersetup{colorlinks=true,
            linkcolor=blue!50!black,
            citecolor=blue!50!black,
            urlcolor=blue!50!black}
\newcommand{\qsixR}{\mathbb{R}}
\newcommand{\qsixE}{\mathbb{E}}
\newcommand{\qsixPr}{\mathbb{P}}
\newcommand{\qsixeps}{\varepsilon}
\newcommand{\qsixTr}{\operatorname{Tr}}
\newcommand{\qsixTrs}{\operatorname{Tr}_{\sigma}}
\newcommand{\qsixdiag}{\operatorname{diag}}
\newcommand{\qsixopnorm}[1]{\left\lVert #1 \right\rVert}
\newcommand{\qsixdel}{\delta}
\newcommand{\qsixSpan}{\operatorname{span}}
\newcommand{\qsixim}{\operatorname{im}}
\newcommand{\qsixrank}{\operatorname{rank}}
\newcommand{\qsixtL}{\widetilde{L}}
\newcommand{\qsixLt}{L^{\dagger}}
\newcommand{\qsixLth}{L^{\dagger/2}}
\newcommand{\qsixBarrier}[2]{\mathcal{B}^{#1}_{#2}}
\newtheorem{qsixtheorem}{Theorem}
\newtheorem{qsixlemma}[qsixtheorem]{Lemma}
\newtheorem{qsixproposition}[qsixtheorem]{Proposition}
\newtheorem{qsixclaim}[qsixtheorem]{Claim}
\title{$\qsixeps$-Light Subsets in Graphs}
\author{}
\date{}
\maketitle


\section*{Statement and Answer}

Let $G=(V,E,w)$ be a weighted graph on $n=|V|$ vertices with Laplacian
$L=\sum_{(s,t)\in E}w(s,t)\,(\delta_s-\delta_t)(\delta_s-\delta_t)^{\!\top}$.
For $S\subseteq V$ let $G_S=(V,E(S,S))$ keep only edges with both endpoints in
$S$ and let $L_S$ be its Laplacian (same ambient dimension as $L$). $S$ is
\emph{$\qsixeps$-light} iff $\qsixeps L - L_S\succeq 0$.

\medskip
\noindent\textbf{Theorem.} \emph{For every weighted graph $G=(V,E,w)$ and every
$\qsixeps\in(0,1)$, there exists an $\qsixeps$-light subset $S\subseteq V$ with
$|S|\ge \lfloor \qsixeps n/42\rfloor$.}

\medskip
\noindent The answer is \textbf{YES}, with the explicit constant $c=1/42$.

\section*{Notation and basic facts}

Write $\qsixLt$ for the Moore--Penrose pseudoinverse and $\qsixLth$ for its PSD square
root. Define
\[
  \qsixtL_S      := \qsixLth L_S \qsixLth, \qquad
  \qsixtL_{S,T}  := \qsixLth L_{S,T}  \qsixLth, \qquad
  \qsixtL_{S,t}  := \qsixLth L_{S,\{t\}}\qsixLth.
\]
Then $\qsixLth L \qsixLth = \Pi$, the orthogonal projector onto $\qsixim L$
(the column space $\qsixSpan\{L\}$). We record the basic kernel containment and
the resulting equivalence, which is the foundational reduction used throughout.

\begin{qsixlemma}[Kernel containment and operator-norm reduction]\label{lem:opnorm-form}
For every $S\subseteq V$ we have $\ker L\subseteq\ker L_S$. Consequently
$\qsixtL_S=\qsixLth L_S\qsixLth$ is supported on $\qsixim L$, i.e.
$\Pi\qsixtL_S=\qsixtL_S\Pi=\qsixtL_S$, and for every $\qsixeps\ge 0$,
\begin{equation}\label{eq:opnorm-form}
  \qsixeps L\succeq L_S
  \;\Longleftrightarrow\;
  \qsixeps\Pi\succeq\qsixtL_S
  \;\Longleftrightarrow\;
  \qsixopnorm{\qsixtL_S}_{\mathrm{op}} \;\le\; \qsixeps.
\end{equation}
\end{qsixlemma}

\begin{proof}
\emph{Kernel containment.} The Laplacian $L$ of $G$ is PSD, so
$\ker L=\ker L^{1/2}=\{x:x^\top L x=0\}$. For a graph Laplacian,
$x^\top L x=\sum_{(s,t)\in E}w(s,t)(x_s-x_t)^2$, which vanishes iff $x$ is constant
on every connected component of $G$; hence $\ker L=\qsixSpan\{\mathbf 1_C:C\text{ a
component of }G\}$. Now let $\{a,b\}$ be any edge of $G_S$ (so $a,b\in S$ and
$\{a,b\}\in E$); its two endpoints lie in the same component $C$ of $G$, on which
each kernel vector $x\in\ker L$ is constant, so $x_a-x_b=0$. Therefore
\[
  x^\top L_S x=\sum_{(s,t)\in E(S,S)}w(s,t)(x_s-x_t)^2=0
  \qquad\text{for all }x\in\ker L,
\]
and since $L_S\succeq 0$ this gives $L_S x=0$, i.e. $\ker L\subseteq\ker L_S$.
Equivalently $\qsixim L_S=(\ker L_S)^\perp\subseteq(\ker L)^\perp=\qsixim L$.

\emph{Support of $\qsixtL_S$.} Since $\Pi$ is the orthogonal projector onto
$\qsixim L=(\ker L)^\perp$ and $\qsixLth$ has range $\qsixim L$ and kernel $\ker L$,
we have $\Pi\qsixLth=\qsixLth\Pi=\qsixLth$; hence
$\Pi\qsixtL_S=\Pi\qsixLth L_S\qsixLth=\qsixLth L_S\qsixLth=\qsixtL_S$ and similarly
$\qsixtL_S\Pi=\qsixtL_S$. Thus $\qsixtL_S=\Pi\qsixtL_S\Pi$ is supported on
$\qsixim L$.

\emph{Equivalence.} Decompose $\qsixR^V=\qsixim L\oplus\ker L$ orthogonally and
write any $x=x_1+x_0$ with $x_1\in\qsixim L$, $x_0\in\ker L$. Because
$L x_0=0$ and (just proved) $L_S x_0=0$, the matrix $A:=\qsixeps L-L_S$ satisfies
$A x_0=0$, so $A$ is supported on $\qsixim L$:
$x^\top A x=x_1^\top A x_1$ for all $x$. Hence $A\succeq 0$ iff $A\succeq 0$ on
$\qsixim L$, i.e. iff $y^\top A y\ge 0$ for all $y\in\qsixim L$.
On $\qsixim L$ the map $y\mapsto \qsixLth y$ is a linear bijection of $\qsixim L$
onto itself with inverse $L^{1/2}$ (since $\qsixLth L^{1/2}=L^{1/2}\qsixLth=\Pi$
restricts to the identity on $\qsixim L$). For $y\in\qsixim L$ set $z:=L^{1/2}y\in\qsixim L$;
then $y=\qsixLth z$ and
\[
  y^\top A y
  =z^\top \qsixLth(\qsixeps L-L_S)\qsixLth z
  =z^\top(\qsixeps\Pi-\qsixtL_S)z ,
\]
using $\qsixLth L\qsixLth=\Pi$. As $y$ ranges over $\qsixim L$ so does $z$, and on
$\ker L$ both $\Pi$ and $\qsixtL_S$ vanish, so the condition ``$y^\top Ay\ge0$ for
all $y\in\qsixim L$'' is equivalent to $\qsixeps\Pi\succeq\qsixtL_S$ on all of
$\qsixR^V$. This proves $\qsixeps L\succeq L_S\iff\qsixeps\Pi\succeq\qsixtL_S$.
Finally, since $\qsixtL_S\succeq0$ is supported on $\qsixim L$ and $\Pi$ is the
identity there, $\qsixeps\Pi\succeq\qsixtL_S$ holds iff every eigenvalue of
$\qsixtL_S$ is at most $\qsixeps$, i.e. iff
$\qsixopnorm{\qsixtL_S}_{\mathrm{op}}\le\qsixeps$.
\end{proof}

Reduction~\eqref{eq:opnorm-form} is what lets us pass between the semidefinite
condition $\qsixeps L\succeq L_S$ and the operator-norm bound on $\qsixtL_S$ in all
that follows.

\medskip
\noindent\emph{Leverage scores.}\quad The leverage score of an edge $(s,t)$ is
\[
  \ell(s,t)
   := w(s,t)\,(\delta_s-\delta_t)^{\!\top}\qsixLt(\delta_s-\delta_t)
    = \qsixTr\!\bigl(L_{\{s\},\{t\}}\qsixLt\bigr),
\]
and zero on non-edges. For $S,T\subseteq V$ set
$\ell(S,T):=\sum_{s\in S,t\in T}\ell(s,t)$ and write
$\ell(S):=\ell(S,V\setminus S)$ for the boundary leverage. By linearity
$\ell(S,T)=\qsixTr(\qsixtL_{S,T})$. We record the following general form of
Foster's identity, valid for an arbitrary (possibly disconnected) graph.

\begin{qsixproposition}[Foster's identity]\label{prop:foster}
Let $G$ have $c$ connected components. Then
\[
  \sum_{(s,t)\in E}\ell(s,t)\;=\;\qsixrank L\;=\;n-c\;\le\;n-1 .
\]
\end{qsixproposition}

\begin{proof}
Summing the edge Laplacians recovers the full Laplacian,
$\sum_{(s,t)\in E}w(s,t)(\delta_s-\delta_t)(\delta_s-\delta_t)^{\!\top}=L$, so by
linearity of the trace and the definition of $\ell$,
\[
  \sum_{(s,t)\in E}\ell(s,t)
   =\sum_{(s,t)\in E}\qsixTr\!\bigl(w(s,t)(\delta_s-\delta_t)(\delta_s-\delta_t)^{\!\top}\qsixLt\bigr)
   =\qsixTr\!\bigl(L\qsixLt\bigr).
\]
Now $L\qsixLt=\Pi$ is the orthogonal projector onto $\qsixim L$, whose trace equals
$\dim\qsixim L=\qsixrank L$; hence $\sum_{(s,t)\in E}\ell(s,t)=\qsixrank L$.
Finally, as computed in the proof of Lemma~\ref{lem:opnorm-form},
$\ker L=\qsixSpan\{\mathbf 1_C:C\text{ a component of }G\}$, and the indicator
vectors $\mathbf 1_C$ of the $c$ distinct components are linearly independent
(their supports are disjoint and nonempty), so $\dim\ker L=c$ and
$\qsixrank L=n-c$. Since $G$ has at least one vertex it has $c\ge 1$ component,
giving $n-c\le n-1$. (For connected $G$, $c=1$ and the bound is the equality
$\sum_{(s,t)\in E}\ell(s,t)=n-1$.)
\end{proof}

\medskip
\noindent\emph{Bounded-rank upper barrier.}\quad For a symmetric $A$ with
eigenvalues $\lambda_1\ge\cdots\ge\lambda_n$ and a real $u$, define
\[
  \qsixBarrier{u}{\sigma}(A) \;:=\;
  \begin{cases}
    \displaystyle\sum_{i=1}^{\sigma}\dfrac{1}{u-\lambda_i} & \text{if } u>\lambda_1,\\[2pt]
    \infty & \text{otherwise,}
  \end{cases}
\]
and the top-$\sigma$ trace $\qsixTrs(A):=\sum_{i=1}^{\sigma}\lambda_i$. We write
$\qsixBarrier{u}{\sigma}(S)$ for $\qsixBarrier{u}{\sigma}(\qsixtL_S)$, and note
$\qsixBarrier{u}{\sigma}(A)=\qsixTrs\bigl((uI-A)^{-1}\bigr)$ when $u>\lambda_1$.

For the rest of this note we fix
\[
  \qsixdel := \frac{21}{n}, \qquad
  \phi := \frac{n}{21}, \qquad
  \sigma := \lfloor \qsixeps n/42 \rfloor.
\]

\section*{Proof}

We build $S$ greedily, adding one vertex at a time and increasing the barrier
parameter $u$ by $\qsixdel$ at each step. Throughout, two invariants are
maintained: the bounded-rank barrier stays at most $\phi$, and the boundary
leverage stays at most $4|S|$.

\subsection*{Step 1: Initial set and target}

By Foster's identity (Proposition~\ref{prop:foster}),
$\sum_{(s,t)\in E}\ell(s,t)=n-c\le n-1$, where $c\ge1$ is the number of connected
components of $G$; this bound holds for every graph, connected or not. Therefore
\[
  \sum_{v\in V}\ell(\{v\},V\setminus\{v\})
   \;=\; 2\!\sum_{(s,t)\in E}\ell(s,t)
   \;\le\; 2(n-1).
\]
The average boundary leverage is less than $2$. By Markov's inequality with
threshold $4$, at most $2(n-1)/4 < n/2$ vertices satisfy
$\ell(\{v\},V\setminus\{v\})>4$; hence more than half of all vertices have
$\ell(\{v\},V\setminus\{v\})\le 4$. Pick any such $v_0$ and set
$S_0:=\{v_0\}$, $u_0:=\qsixeps/2$.

Then $\ell(S_0)=\ell(\{v_0\},V\setminus\{v_0\})\le 4=4|S_0|$
(leverage invariant at step $0$).
Since $G_{S_0}$ has no edges, $\qsixtL_{S_0}=0$, and
\[
  \qsixBarrier{u_0}{\sigma}(S_0) \;=\; \frac{\sigma}{u_0}
   \;\le\; \frac{\qsixeps n/42}{\qsixeps/2} \;=\; \frac{n}{21} \;=\; \phi
\]
(barrier invariant at step $0$).

\begin{qsixclaim}[Rank bound for bounded-rank barrier]\label{cl:rank-barrier}
If $|S|=\sigma+1$, then $\qsixrank(L_S) \le \sigma$, and thus $\qsixrank(\qsixtL_S) \le \sigma$.
Consequently, $\qsixtL_S$ has at most $\sigma$ nonzero eigenvalues. Therefore, if
$\qsixBarrier{u_\sigma}{\sigma}(S)<\infty$, then every nonzero eigenvalue of $\qsixtL_S$
lies strictly below $u_\sigma$, so $\qsixopnorm{\qsixtL_S}_{\mathrm{op}}\le u_\sigma$.
\end{qsixclaim}

\begin{proof}
The Laplacian $L_S$ of a graph on $|S|$ vertices has rank at most $|S|-1$ (for a
connected component, the rank equals the number of vertices minus $1$; for
multiple components, it is the number of vertices minus the number of components,
which is at most $|S|-1$). Thus $\qsixrank(L_S)\le|S|-1=\sigma$.

By construction, $\qsixtL_S=\qsixLth L_S\qsixLth$, and matrix multiplication
preserves rank order, so $\qsixrank(\qsixtL_S)\le\qsixrank(L_S)\le\sigma$.

A PSD matrix with rank at most $\sigma$ has exactly $\sigma - (\sigma - \text{rank})
= \text{rank}$ nonzero eigenvalues, so $\qsixtL_S$ has at most $\sigma$ nonzero
eigenvalues. When $\qsixBarrier{u_\sigma}{\sigma}(S)<\infty$, the top $\sigma$
eigenvalues $\lambda_1\ge\cdots\ge\lambda_\sigma$ of $\qsixtL_S$ all satisfy
$\lambda_i < u_\sigma$. Since there are at most $\sigma$ nonzero eigenvalues in
total, these are all the nonzero eigenvalues, and the remaining eigenvalues are
zero. Thus $\qsixopnorm{\qsixtL_S}_{\mathrm{op}} = \lambda_1 < u_\sigma$.
\end{proof}

After $\sigma$ greedy steps, the barrier will sit at $u_\sigma=u_0+\sigma\qsixdel$.
We chose parameters so that
\[
  u_\sigma = \tfrac{\qsixeps}{2} + \sigma\cdot\tfrac{21}{n}
           \le \tfrac{\qsixeps}{2} + \tfrac{\qsixeps}{2} = \qsixeps,
\]
so once we reach $|S|=\sigma+1$ with $\qsixBarrier{u_\sigma}{\sigma}(S)$ finite, by
Claim~\ref{cl:rank-barrier} every eigenvalue of $\qsixtL_S$ lies below $\qsixeps$ and
\eqref{eq:opnorm-form} gives $\qsixeps L\succeq L_S$.

\subsection*{Step 2: Greedy lemma}

\begin{qsixlemma}\label{lem:greedy}
Suppose $|S|\le\sigma$, $\ell(S)\le 4|S|$, and $\qsixBarrier{u}{\sigma}(S)\le\phi$. Then
there exists $t\notin S$ with
\[
  \qsixBarrier{u+\qsixdel}{\sigma}(S\cup\{t\})\le\phi
  \qquad\text{and}\qquad
  \ell(S\cup\{t\})\le\ell(S)+4.
\]
\end{qsixlemma}

\begin{proof}[Proof of Lemma~\ref{lem:greedy}, modulo the two averaging bounds]
Set $m:=n-|S|=|V\setminus S|\ge n-\sigma>0$ and define the two candidate sets
\[
  \mathcal{A}:=\{t\notin S:\ \ell(t,V\setminus S)\le 4\},
  \qquad
  \mathcal{B}:=\{t\notin S:\ U(S,t)<1\},
\]
where $U(S,t)$ is the quantity defined in \eqref{eq:U}. In Step~3 we prove
$|\mathcal{A}|>m/2$ and in Step~4 we prove $|\mathcal{B}|>m/2$. Granting these,
\[
  |\mathcal{A}|+|\mathcal{B}|>m=|V\setminus S|,
\]
so $\mathcal{A}\cap\mathcal{B}\ne\varnothing$; fix any
$t\in\mathcal{A}\cap\mathcal{B}$. We verify the two conclusions of the lemma
for this $t$.

\medskip\noindent\emph{(i) Leverage bound.}\quad Since adding the single
vertex $t$ to $S$ creates no new edges inside $\{t\}$, the only new edges of
$G_{S\cup\{t\}}$ relative to $G_S$ are the cross-edges between $t$ and $S$.
Hence for $t\notin S$,
\[
  \ell(S\cup\{t\})
  =\ell(S\cup\{t\},V\setminus(S\cup\{t\}))
  \le\ell(S\cup\{t\},V\setminus S)
  =\ell(S,V\setminus S)+\ell(t,V\setminus S)
  =\ell(S)+\ell(t,V\setminus S),
\]
where the inequality drops the (non-negative) boundary terms toward the
removed vertex $t$, and the third equality splits the boundary of $S\cup\{t\}$
into the part emanating from $S$ and the part emanating from $t$. Because
$t\in\mathcal{A}$ we have $\ell(t,V\setminus S)\le 4$, so
\[
  \ell(S\cup\{t\})\le\ell(S)+4,
\]
which is the second conclusion.

\medskip\noindent\emph{(ii) Barrier bound.}\quad We apply
Lemma~\ref{lem:bu} with the substitution
\[
  A:=\qsixtL_S,\qquad B:=\qsixtL_{S,t},\qquad
  M=(u+\qsixdel)I-A=(u+\qsixdel)I-\qsixtL_S,
\]
matching the matrix $M$ of \eqref{eq:U}. We check the hypotheses.

\emph{$B\succeq0$ and the additive structure.} $\qsixtL_{S,t}=\qsixLth
L_{S,\{t\}}\qsixLth$ is a conjugate of the PSD matrix $L_{S,\{t\}}$
(a sum of edge Laplacians $w(s,t)(\delta_s-\delta_t)(\delta_s-\delta_t)^{\!\top}$
over $s\in S$), hence $B\succeq0$. Moreover $G_{S\cup\{t\}}$ decomposes as the
edges inside $S$ plus the cross-edges from $t$ to $S$ (no edge lies inside
$\{t\}$), so $L_{S\cup\{t\}}=L_S+L_{S,\{t\}}$ and therefore
\[
  \qsixtL_{S\cup\{t\}}=\qsixtL_S+\qsixtL_{S,t}=A+B .
\]
Thus the conclusion $\qsixBarrier{u+\qsixdel}{\sigma}(A+B)\le\qsixBarrier{u}{\sigma}(A)$
of Lemma~\ref{lem:bu} is exactly
$\qsixBarrier{u+\qsixdel}{\sigma}(S\cup\{t\})\le\qsixBarrier{u}{\sigma}(S)$.

\emph{Rank condition $\qsixrank(B)\le\sigma$.} Since $|S|\le\sigma$, the graph
$G_{S\cup\{t\}}$ has at most $\sigma+1$ non-isolated vertices, so
$\qsixrank(L_{S,\{t\}})\le\qsixrank(L_{S\cup\{t\}})\le(\sigma+1)-1=\sigma$
by the Laplacian rank bound; conjugation by $\qsixLth$ does not increase rank,
so $\qsixrank(B)=\qsixrank(\qsixtL_{S,t})\le\sigma$. (Directly:
$L_{S,\{t\}}=\sum_{s\in S}w(s,t)(\delta_s-\delta_t)(\delta_s-\delta_t)^{\!\top}$
is a sum of at most $|S|\le\sigma$ rank-one terms.)

\emph{Finiteness $\qsixBarrier{u}{\sigma}(A)<\infty$.} This is the hypothesis
$\qsixBarrier{u}{\sigma}(S)\le\phi<\infty$ of Lemma~\ref{lem:greedy}.

\emph{The scalar hypothesis.} With these substitutions, the left-hand side of
the hypothesis of Lemma~\ref{lem:bu},
\[
  \frac{\qsixTrs(M^{-2}B)}{\qsixBarrier{u}{\sigma}(A)-\qsixBarrier{u+\qsixdel}{\sigma}(A)}
   + \qsixTrs(M^{-1}B),
\]
is \emph{verbatim} the quantity $U(S,t)$ of \eqref{eq:U}
(here $\qsixBarrier{u}{\sigma}(A)=\qsixBarrier{u}{\sigma}(S)$ and
$\qsixBarrier{u+\qsixdel}{\sigma}(A)=\qsixBarrier{u+\qsixdel}{\sigma}(S)$ because
$A=\qsixtL_S$). Since $t\in\mathcal{B}$ means $U(S,t)<1$, the hypothesis of
Lemma~\ref{lem:bu} holds. The lemma therefore yields
\[
  \qsixBarrier{u+\qsixdel}{\sigma}(S\cup\{t\})
   =\qsixBarrier{u+\qsixdel}{\sigma}(A+B)
   \le\qsixBarrier{u}{\sigma}(A)
   =\qsixBarrier{u}{\sigma}(S)\le\phi,
\]
which is the first conclusion. This completes the proof of
Lemma~\ref{lem:greedy}, contingent on the bounds $|\mathcal{A}|>m/2$
(Step~3) and $|\mathcal{B}|>m/2$ (Step~4) established below.
\end{proof}

\medskip
\noindent The two deferred bounds are proved by averaging: more than half the
candidates $t$ keep the leverage budget (Step~3, giving $|\mathcal{A}|>m/2$),
and more than half keep the barrier (Step~4, giving $|\mathcal{B}|>m/2$).

\subsection*{Step 3: Leverage-score averaging}

\begin{qsixclaim}\label{cl:Tlev}
For every $T\subseteq V$, $\sum_{t\in T}\ell(t,T)\le 2(|T|-1)$.
\end{qsixclaim}
\begin{proof}
Each undirected edge $\{s,t\}\subseteq T$ contributes to both $\ell(s,T)$ and
$\ell(t,T)$, so $\sum_{t\in T}\ell(t,T)=2\sum_{\{s,t\}\subseteq T}\ell(s,t)
=2\qsixTr(L_T\qsixLt)$.

Since $\qsixLt=\qsixLth\qsixLth$, cyclicity of trace gives
\[
  \qsixTr(L_T\qsixLt)
  =\qsixTr(L_T\cdot\qsixLth\cdot\qsixLth)
  =\qsixTr(\qsixLth\cdot L_T\cdot\qsixLth)
  =\qsixTr(\qsixtL_T),
\]
where $\qsixtL_T:=\qsixLth L_T\qsixLth\succeq 0$ is symmetric PSD (using the same
notation as in the Notation section with $T$ in place of $S$). Since $L_T\preceq L$,
conjugating by the PSD matrix $\qsixLth$ preserves the order:
$\qsixtL_T\preceq\qsixLth L\qsixLth=\Pi$, so all eigenvalues of $\qsixtL_T$ lie in
$[0,1]$. Since $\qsixrank(\qsixtL_T)\le\qsixrank(L_T)\le|T|-1$, the matrix $\qsixtL_T$
has at most $|T|-1$ nonzero eigenvalues, each in $[0,1]$, giving
$\qsixTr(\qsixtL_T)\le|T|-1$.
\end{proof}

We bound $|\mathcal{A}|$. Claim~\ref{cl:Tlev} applied to $T=V\setminus S$ gives
\[
  \sum_{t\notin S}\ell(t,V\setminus S)\le 2\bigl(|V\setminus S|-1\bigr)<2m,
\]
so the average of $\ell(t,V\setminus S)$ over the $m$ candidates $t\notin S$ is
less than $2$. Since $\ell(t,V\setminus S)\ge0$, Markov's inequality at
threshold $4$ bounds the number of $t\notin S$ with $\ell(t,V\setminus S)>4$ by
$\tfrac{2m}{4}=\tfrac{m}{2}$; equivalently fewer than $m/2$ candidates fail the
condition. Hence
\[
  |\mathcal{A}|=\#\{t\notin S:\ \ell(t,V\setminus S)\le4\}>m-\tfrac{m}{2}=\tfrac{m}{2},
\]
which is the bound used in the proof of Lemma~\ref{lem:greedy}. (Note the
boundary-leverage inequality $\ell(S\cup\{t\})\le\ell(S)+\ell(t,V\setminus S)$
used there was established in part (i) of that proof.)

\subsection*{Step 4: Barrier averaging}

This is the technical heart. Here we prove the bound $|\mathcal{B}|>m/2$ used
in the proof of Lemma~\ref{lem:greedy}, where
$\mathcal{B}=\{t\notin S:\ U(S,t)<1\}$ and
\begin{equation}\label{eq:U}
  U(S,t) \;:=\;
  \frac{\qsixTrs(M^{-2}\qsixtL_{S,t})}{\qsixBarrier{u}{\sigma}(S)-\qsixBarrier{u+\qsixdel}{\sigma}(S)}
   + \qsixTrs(M^{-1}\qsixtL_{S,t}) \;<\; 1,
\end{equation}
with $M:=(u+\qsixdel)I-\qsixtL_S$. The mixed-product traces
$\qsixTrs(M^{-2}\qsixtL_{S,t})$ and $\qsixTrs(M^{-1}\qsixtL_{S,t})$ in \eqref{eq:U}
are read in the symmetrized sense of the convention stated with
Lemma~\ref{lem:bu}: $\qsixTrs(A'N):=\qsixTrs(N^{1/2}A'N^{1/2})$ for symmetric
$A'$ and PSD $N$ (here $M^{-2},M^{-1}$ are symmetric and $\qsixtL_{S,t}\succeq0$);
the same convention is recalled and applied in the $\Pi_S^\perp$/$\Pi_S$
computation below. We emphasize that the expression \eqref{eq:U}
is, term by term, the left-hand side of the hypothesis of Lemma~\ref{lem:bu}
under the substitution $A=\qsixtL_S$, $B=\qsixtL_{S,t}$; this is what makes the
event $U(S,t)<1$ the precise admissibility condition for adding $t$, as used in
part (ii) of the proof of Lemma~\ref{lem:greedy}. (The rank condition
$\qsixrank(\qsixtL_{S,t})\le\sigma$ required by Lemma~\ref{lem:bu} was also
checked there, using $|S|\le\sigma$.) We claim
\begin{equation}\label{eq:total}
  \sum_{t\notin S} U(S,t) \;\le\; \tfrac{5}{\qsixdel} + 5\phi.
\end{equation}
Granting \eqref{eq:total}, since each $U(S,t)\ge 0$ the average of $U(S,t)$ over
the $m=n-|S|$ candidates is at most $\tfrac1m(\tfrac5\qsixdel+5\phi)$, so by
Markov's inequality at threshold $1$ the number of $t\notin S$ with
$U(S,t)\ge1$ is at most
\[
  \frac{1}{m}\Bigl(\tfrac{5}{\qsixdel}+5\phi\Bigr)
  =\frac{1}{n-|S|}\Bigl(\tfrac{5n}{21}+\tfrac{5n}{21}\Bigr)
  \le\frac{1}{41n/42}\cdot\tfrac{10n}{21}
  =\frac{42}{41}\cdot\tfrac{10}{21}
  =\frac{20}{41}<\frac12 ,
\]
using $\qsixdel=21/n$, $\phi=n/21$ and $|S|\le\sigma\le n/42$ (so
$m=n-|S|\ge 41n/42$). Therefore the number of $t\notin S$ with $U(S,t)<1$
exceeds $m-\tfrac{m}{2}=\tfrac{m}{2}$, i.e.
\[
  |\mathcal{B}|=\#\{t\notin S:\ U(S,t)<1\}>\tfrac{m}{2},
\]
as required.

\medskip
\noindent\emph{Proof of \eqref{eq:total}.}\quad Let $\Pi_S$ project onto
$\qsixim\qsixtL_S$, and $\Pi_S^\perp:=I-\Pi_S$. Throughout this proof $p\in\{-1,-2\}$.

\emph{The matrix $M$ is invertible, so $M^p$ ($p<0$) is the ordinary inverse
power.} Indeed $M=(u+\qsixdel)I-\qsixtL_S$ is symmetric with eigenvalues
$u+\qsixdel-\lambda_i(\qsixtL_S)$. The hypothesis $\qsixBarrier{u}{\sigma}(S)<\infty$
forces $u>\lambda_1(\qsixtL_S)$ (recall $\qsixBarrier{u}{\sigma}(A)=\infty$ unless
$u>\lambda_1(A)$), hence $u+\qsixdel>u>\lambda_1(\qsixtL_S)\ge\lambda_i(\qsixtL_S)$
for all $i$. Thus every eigenvalue of $M$ is strictly positive, so $M\succ0$ is
invertible and $M^{-1},M^{-2}$ are well defined symmetric PSD matrices; this is the
only meaning attached to $M^p$ below.

\emph{Block decomposition of $M^p$.} The projector $\Pi_S$ onto $\qsixim\qsixtL_S$
is a spectral projector of $\qsixtL_S$, so $\qsixtL_S\Pi_S=\Pi_S\qsixtL_S=\qsixtL_S$
and consequently $M\Pi_S=\Pi_S M$. Therefore $M$ (and hence any power of $M$)
leaves both $\qsixim\Pi_S$ and $\qsixim\Pi_S^\perp$ invariant, and $M^p$ commutes
with $\Pi_S$. On $\qsixim\Pi_S^\perp=\ker\qsixtL_S$ we have $\qsixtL_S=0$, so $M$
acts there as $(u+\qsixdel)I$ and $M^p$ acts as $(u+\qsixdel)^pI$. Splitting
$I=\Pi_S+\Pi_S^\perp$ and using $M^p\Pi_S^\perp=(u+\qsixdel)^p\Pi_S^\perp$ together
with $\Pi_S M^p\Pi_S=M^p\Pi_S$ gives the exact identity
\begin{equation}\label{eq:Mp-block}
  M^p \;=\; M^p\Pi_S^\perp + M^p\Pi_S
        \;=\; (u+\qsixdel)^p\,\Pi_S^\perp \;+\; \Pi_S M^p\Pi_S,
  \qquad p\in\{-1,-2\},
\end{equation}
in which every term is an ordinary (well-defined) matrix; the second summand
$\Pi_S M^p\Pi_S$ is the symmetric PSD operator agreeing with $M^p$ on
$\qsixim\Pi_S$ and vanishing on $\qsixim\Pi_S^\perp$. (Equivalently, on
$\qsixim\Pi_S$ this is the $p$-th power of the positive-definite restriction
$M|_{\qsixim\Pi_S}=(u+\qsixdel)\Pi_S-\qsixtL_S$, but we never form a negative power
of the ambient singular matrix $(u+\qsixdel)\Pi_S-\qsixtL_S$ itself.)

\emph{The symmetrized trace convention.} Since $M^p$ is symmetric and
$\qsixtL_{S,t}$ is PSD but their product is not necessarily symmetric, we recall
the symmetrized convention stated with Lemma~\ref{lem:bu}: for symmetric $A$ and
PSD $N$,
\[
  \qsixTrs(AN) \;:=\; \qsixTrs(N^{1/2}AN^{1/2}),
\]
the right-hand side summing the top-$\sigma$ eigenvalues of the symmetric matrix
$N^{1/2}AN^{1/2}$.  (Cyclicity of trace gives $\qsixTr(AN)=\qsixTr(N^{1/2}AN^{1/2})$,
so this is consistent with the ordinary trace.)  Applying this with $N=\qsixtL_{S,t}$
and substituting the decomposition \eqref{eq:Mp-block},
\begin{align*}
\qsixTrs(M^p\qsixtL_{S,t})
  &:= \qsixTrs\!\bigl(\qsixtL_{S,t}^{1/2}M^p\qsixtL_{S,t}^{1/2}\bigr) \\
  &= \qsixTrs\!\bigl((u+\qsixdel)^p\qsixtL_{S,t}^{1/2}\Pi_S^\perp\qsixtL_{S,t}^{1/2}
     + \qsixtL_{S,t}^{1/2}(\Pi_S M^p\Pi_S)\qsixtL_{S,t}^{1/2}\bigr).
\end{align*}
Both summands are symmetric PSD (each is the conjugation of a symmetric PSD matrix
by the PSD factor $\qsixtL_{S,t}^{1/2}$: indeed $\Pi_S^\perp$ is a projector,
$(u+\qsixdel)^p>0$, and $\Pi_S M^p\Pi_S\succeq0$ since $M^p\succeq0$), so Ky Fan
subadditivity of $\qsixTrs$ \cite[Thm.~4.3.47a]{HJ12} applies:
\begin{equation}\label{eq:split}
  \qsixTrs(M^p\qsixtL_{S,t}) \;\le\;
  \qsixTrs\!\bigl((u+\qsixdel)^p\Pi_S^\perp\qsixtL_{S,t}\bigr)
  + \qsixTrs\!\bigl((\Pi_S M^p\Pi_S)\,\qsixtL_{S,t}\bigr).
\end{equation}

\smallskip
\emph{$\Pi_S^\perp$ piece.}\quad
Since $\Pi_S^\perp\preceq I$, conjugating by $\qsixtL_{S,t}^{1/2}\succeq0$ preserves
the semidefinite order: $\qsixtL_{S,t}^{1/2}\Pi_S^\perp\qsixtL_{S,t}^{1/2}\preceq
\qsixtL_{S,t}^{1/2}I\qsixtL_{S,t}^{1/2}=\qsixtL_{S,t}$.
Taking the trace (which is monotone for PSD matrices: $0\preceq P\preceq Q\Rightarrow\qsixTr(P)\le\qsixTr(Q)$) and using cyclicity ($\qsixTr(\qsixtL_{S,t}^{1/2}\Pi_S^\perp\qsixtL_{S,t}^{1/2})=\qsixTr(\Pi_S^\perp\qsixtL_{S,t})$) gives
$\qsixTr(\Pi_S^\perp\qsixtL_{S,t})\le\qsixTr(\qsixtL_{S,t})=\ell(S,t)$.
Combining this with the symmetrized convention and $\qsixTrs(B)\le\qsixTr(B)$ for PSD $B$:
\[
  \qsixTrs(\Pi_S^\perp\qsixtL_{S,t})
    =\qsixTrs(\qsixtL_{S,t}^{1/2}\Pi_S^\perp\qsixtL_{S,t}^{1/2})
    \le\qsixTr(\qsixtL_{S,t}^{1/2}\Pi_S^\perp\qsixtL_{S,t}^{1/2})
    =\qsixTr(\Pi_S^\perp\qsixtL_{S,t})
    \le\qsixTr(\qsixtL_{S,t})=\ell(S,t),
\]
and summing,
\[
  \sum_{t\notin S}\qsixTrs\!\bigl((u+\qsixdel)^p\Pi_S^\perp\qsixtL_{S,t}\bigr)
   \le (u+\qsixdel)^p\,\ell(S)\le (u+\qsixdel)^p\cdot 4|S|.
\]

\smallskip
\emph{$\Pi_S$ piece.}\quad
We first justify $\qsixtL_{S,V\setminus S}\preceq\Pi$.
Partition the edge set $E$ into edges with both endpoints in $S$,
edges with both endpoints in $V\setminus S$, and cross-edges:
\[
  L \;=\; L_S \;+\; L_{V\setminus S} \;+\; L_{S,V\setminus S},
\]
where $L_{V\setminus S}$ denotes the Laplacian of the induced subgraph
on $V\setminus S$ (lifted to the ambient space $\qsixR^V$).  Both $L_S$
and $L_{V\setminus S}$ are Laplacians of non-negatively weighted graphs,
hence PSD.  Therefore
\[
  L - L_{S,V\setminus S} \;=\; L_S + L_{V\setminus S} \;\succeq\; 0,
\]
i.e., $L_{S,V\setminus S}\preceq L$.  Conjugating by the PSD matrix
$\qsixLth$ preserves the semidefinite order: for any $x\in\qsixR^V$,
$x^{\!\top}\qsixLth L_{S,V\setminus S}\qsixLth x
 =(\qsixLth x)^{\!\top}L_{S,V\setminus S}(\qsixLth x)
 \le (\qsixLth x)^{\!\top}L(\qsixLth x)
 =x^{\!\top}\qsixLth L\qsixLth x$, so
\[
  \qsixtL_{S,V\setminus S}
   \;=\; \qsixLth L_{S,V\setminus S}\qsixLth
   \;\preceq\; \qsixLth L\qsixLth \;=\; \Pi.
\]
Hence $\sum_{t\notin S}\qsixtL_{S,t}=\qsixtL_{S,V\setminus S}\preceq\Pi$.
We also note $\Pi_S M^p\Pi_S$ is symmetric PSD (it is $M^p\succeq0$ compressed by
the projector $\Pi_S$) and supported on $\qsixim\Pi_S$.

Denoting $A:=\Pi_S M^p\Pi_S$ (symmetric PSD, supported on $\qsixim\Pi_S$) and
applying the symmetrized convention established above,
\[
\qsixTrs(A\qsixtL_{S,t})
  =\qsixTrs(\qsixtL_{S,t}^{1/2}A\qsixtL_{S,t}^{1/2})
  \le\qsixTr(\qsixtL_{S,t}^{1/2}A\qsixtL_{S,t}^{1/2})
  =\qsixTr(A\qsixtL_{S,t})
\]
($\qsixTrs(B)\le\qsixTr(B)$ for any PSD $B$, and cyclicity of trace in the last step).
Thus $\qsixTrs(A\qsixtL_{S,t})\le\qsixTr(A\qsixtL_{S,t})$ for each $t$. By linearity of trace,
\[
\begin{aligned}
\sum_{t\notin S}\qsixTrs(A\qsixtL_{S,t})
  &\le\sum_{t\notin S}\qsixTr(A\qsixtL_{S,t})\\
  &= \qsixTr\!\Bigl(A\sum_{t\notin S}\qsixtL_{S,t}\Bigr)\\
  &= \qsixTr(A\,\qsixtL_{S,V\setminus S}).
\end{aligned}
\]
Next, using $\qsixtL_{S,V\setminus S}\preceq\Pi$ and $A\succeq0$: the matrix
$\Pi-\qsixtL_{S,V\setminus S}\succeq0$, so $\qsixTr\!\bigl(A(\Pi-\qsixtL_{S,V\setminus S})\bigr)\ge0$
(trace of a product of two PSD matrices is non-negative), giving
\[
\qsixTr(A\,\qsixtL_{S,V\setminus S})\le\qsixTr(A\,\Pi)=\qsixTr(A\Pi)=\qsixTr(\Pi A\Pi),
\]
where the last step uses cyclicity. Since $A=\Pi_S M^p\Pi_S$ is supported on
$\qsixim\Pi_S\subseteq\qsixim\Pi$ (indeed $\Pi\Pi_S=\Pi_S\Pi=\Pi_S$ because
$\qsixim\qsixtL_S\subseteq\qsixim L=\qsixim\Pi$ by Lemma~\ref{lem:opnorm-form}),
we have $\Pi A\Pi=A=\Pi_S M^p\Pi_S$.

We now bound $\qsixTr(\Pi_S M^p\Pi_S)$ by the top-$\sigma$ trace of $M^p$ via the
Ky Fan \emph{maximum principle} (the variational characterization of the sum of the
largest eigenvalues; see Bhatia \cite[Cor.~III.4.4 and Problem~III.6.5]{Bhatia97}
or Horn--Johnson \cite[Cor.~4.3.39]{HJ12}). This is a \emph{distinct} statement
from the Ky Fan subadditivity inequality \cite[Thm.~4.3.47a]{HJ12} used in
\eqref{eq:split} above. For completeness we give the short self-contained proof,
which is all we need here.

\begin{qsixlemma}[Ky Fan maximum principle, compression form]\label{lem:kyfan-max}
Let $B$ be a symmetric matrix with eigenvalues $\mu_1\ge\mu_2\ge\cdots$, and let
$P$ be an orthogonal projector of rank $r$. Then
$\qsixTr(PBP)\le\sum_{i=1}^{r}\mu_i=\qsixTr_r(B)$. If moreover $B\succeq0$ and
$r\le\sigma$, then $\qsixTr_r(B)\le\qsixTrs(B)$.
\end{qsixlemma}

\begin{proof}
Let $\{q_1,\dots,q_r\}$ be an orthonormal basis of $\qsixim P$, so that
$P=\sum_{i=1}^r q_iq_i^{\!\top}$ and, by cyclicity of the trace,
$\qsixTr(PBP)=\qsixTr(PB)=\sum_{i=1}^{r}q_i^{\!\top}Bq_i$. Write the spectral
decomposition $B=\sum_{k}\mu_k e_ke_k^{\!\top}$ with $\{e_k\}$ orthonormal and
$\mu_1\ge\mu_2\ge\cdots$. Setting $c_{ik}:=(q_i^{\!\top}e_k)^2\ge0$ we get
\[
  \qsixTr(PBP)=\sum_{i=1}^{r}\sum_k \mu_k\,c_{ik}=\sum_k \mu_k\,d_k,
  \qquad d_k:=\sum_{i=1}^{r}c_{ik}.
\]
Since $\{q_i\}$ is orthonormal, $\sum_i c_{ik}=\sum_i(e_k^{\!\top}q_i)^2
=\|Pe_k\|^2\in[0,1]$, i.e. $0\le d_k\le1$; and since $\{e_k\}$ is orthonormal,
$\sum_k d_k=\sum_{i=1}^r\sum_k(q_i^{\!\top}e_k)^2=\sum_{i=1}^r\|q_i\|^2=r$.
Thus $(d_k)$ lies in the polytope $\{0\le d_k\le1,\ \sum_k d_k=r\}$, over which the
linear functional $\sum_k\mu_k d_k$ (with $\mu_k$ nonincreasing) is maximized by
putting weight $1$ on the $r$ largest $\mu_k$: indeed
\[
  \sum_k\mu_k d_k-\sum_{i=1}^{r}\mu_i
   =\sum_{k>r}\mu_k d_k-\sum_{i=1}^{r}\mu_i(1-d_i)
   \le\mu_r\sum_{k>r}d_k-\mu_r\sum_{i=1}^{r}(1-d_i)=0,
\]
the inequality using $\mu_k\le\mu_r$ for $k>r$, $\mu_i\ge\mu_r$ for $i\le r$, and
$d_k\ge0$, $1-d_i\ge0$, while the final equality uses
$\sum_{k>r}d_k=r-\sum_{i\le r}d_i=\sum_{i\le r}(1-d_i)$ (from $\sum_kd_k=r$).
Hence $\qsixTr(PBP)\le\sum_{i=1}^{r}\mu_i=\qsixTr_r(B)$. Finally, if $B\succeq0$ then
all $\mu_i\ge0$, so for $r\le\sigma$ the sum of the top $r$ eigenvalues is at most
the sum of the top $\sigma$ eigenvalues: $\qsixTr_r(B)\le\qsixTrs(B)$.
\end{proof}

We apply Lemma~\ref{lem:kyfan-max} with $B:=M^p\succeq0$ and $P:=\Pi_S$. Its rank is
$\qsixrank(\Pi_S)=\qsixrank(\qsixtL_S)\le |S|-1$ (Laplacian rank bound), and since
$|S|\le\sigma$ throughout the induction we have $r:=\qsixrank(\Pi_S)\le\sigma-1\le\sigma$.
The lemma therefore gives
$\qsixTr(\Pi_S M^p\Pi_S)\le\qsixTr_{r}(M^p)\le\qsixTrs(M^p)$.

\smallskip
Since all eigenvalues of $M$ are at most $u+\qsixdel$, for $p<0$ all eigenvalues of
$M^p$ are at least $(u+\qsixdel)^p$, so $\qsixTrs(M^p)\ge\sigma(u+\qsixdel)^p\ge|S|(u+\qsixdel)^p$.
Combining,
\[
  \sum_{t\notin S}\qsixTrs(M^p\qsixtL_{S,t}) \;\le\; 4|S|(u+\qsixdel)^p + \qsixTrs(M^p)
   \;\le\; 5\,\qsixTrs(M^p).
\]
Now we use (Claim~\ref{cl:phi-diff})
$\qsixBarrier{u}{\sigma}(S)-\qsixBarrier{u+\qsixdel}{\sigma}(S)\ge\qsixdel\,\qsixTrs(M^{-2})$ to bound the
first term of $U(S,t)$, and $\qsixTrs(M^{-1})=\qsixBarrier{u+\qsixdel}{\sigma}(S)\le\qsixBarrier{u}{\sigma}(S)\le\phi$ for the
second:
\[
  \sum_{t\notin S}U(S,t)
   \;\le\; \frac{5\,\qsixTrs(M^{-2})}{\qsixdel\,\qsixTrs(M^{-2})} + 5\,\qsixBarrier{u+\qsixdel}{\sigma}(S)
   \;\le\; \frac{5\,\qsixTrs(M^{-2})}{\qsixdel\,\qsixTrs(M^{-2})} + 5\,\qsixBarrier{u}{\sigma}(S)
   \;\le\; \tfrac{5}{\qsixdel}+5\phi,
\]
which is \eqref{eq:total}.\hfill$\square$

\subsection*{Step 5: Combining and conclusion}

We emphasize at the outset what must be produced: a \emph{single} set $S$ that
\emph{simultaneously} satisfies (i) $\qsixeps L-L_S\succeq0$ and
(ii) $|S|\ge\lfloor\qsixeps n/42\rfloor=\sigma$. The greedy construction below
outputs exactly one chain of sets $S_0\subsetneq S_1\subsetneq\cdots\subsetneq
S_\sigma$; its last element $S:=S_\sigma$ is the set we return, and we verify both
(i) and (ii) for this same $S$. There is no separate ``large set'' and ``light
set'': the size is the number of greedy steps and the spectral bound is the
terminal value of the invariant, and both are read off the one chain.

\begin{qsixproposition}[Greedy chain]\label{prop:chain}
There is a chain of vertex sets
\[
  S_0\subsetneq S_1\subsetneq\cdots\subsetneq S_\sigma,
  \qquad |S_k|=k+1,
\]
together with parameters $u_k:=u_0+k\qsixdel=\tfrac{\qsixeps}{2}+k\qsixdel$, such that for
every $k\in\{0,1,\dots,\sigma\}$ the following three invariants hold for the
\emph{same} set $S_k$:
\begin{equation}\label{eq:chain-inv}
  |S_k|=k+1\le\sigma+1,\qquad
  \ell(S_k)\le 4|S_k|,\qquad
  \qsixBarrier{u_k}{\sigma}(S_k)\le\phi .
\end{equation}
\end{qsixproposition}

\begin{proof}
Induction on $k$. \emph{Base case $k=0$:} take $S_0=\{v_0\}$ with $v_0$ chosen in
Step~1. There $|S_0|=1$, $\ell(S_0)\le4=4|S_0|$, and
$\qsixBarrier{u_0}{\sigma}(S_0)=\sigma/u_0\le\phi$ were all verified in Step~1, so
\eqref{eq:chain-inv} holds at $k=0$.

\emph{Inductive step.} Suppose $0\le k<\sigma$ and \eqref{eq:chain-inv} holds for
$S_k$. Then $|S_k|=k+1\le\sigma$, so the hypotheses of Lemma~\ref{lem:greedy}
(namely $|S_k|\le\sigma$, $\ell(S_k)\le4|S_k|$, and
$\qsixBarrier{u_k}{\sigma}(S_k)\le\phi$) are met with $u=u_k$. The lemma supplies a
vertex $t\notin S_k$; set $S_{k+1}:=S_k\cup\{t\}$, so $S_k\subsetneq S_{k+1}$ and
$|S_{k+1}|=k+2$. The two conclusions of the lemma give
\[
  \qsixBarrier{u_k+\qsixdel}{\sigma}(S_{k+1})\le\phi,
  \qquad
  \ell(S_{k+1})\le\ell(S_k)+4 .
\]
Since $u_{k+1}=u_k+\qsixdel$, the first is exactly
$\qsixBarrier{u_{k+1}}{\sigma}(S_{k+1})\le\phi$. For the leverage invariant, the
inductive hypothesis $\ell(S_k)\le4|S_k|=4(k+1)$ gives
$\ell(S_{k+1})\le4(k+1)+4=4(k+2)=4|S_{k+1}|$. And $|S_{k+1}|=k+2\le\sigma+1$.
Thus \eqref{eq:chain-inv} holds for $S_{k+1}$, completing the induction.
\end{proof}

Now fix $S:=S_\sigma$, the terminal set of the chain in
Proposition~\ref{prop:chain}; this is the single set we return. We verify
properties (i) and (ii) for it.

\medskip\noindent\emph{(ii) Size.}\quad By Proposition~\ref{prop:chain},
$|S|=|S_\sigma|=\sigma+1\ge\sigma=\lfloor\qsixeps n/42\rfloor$. (When $\sigma=0$ the
chain is the single set $S=S_0$ of size $1\ge0$, and the bound is the trivial
$|S|\ge0$.)

\medskip\noindent\emph{(i) Lightness of the very same $S$.}\quad The terminal
invariant \eqref{eq:chain-inv} at $k=\sigma$ reads
$\qsixBarrier{u_\sigma}{\sigma}(S)\le\phi<\infty$, with
\[
  u_\sigma=\tfrac{\qsixeps}{2}+\sigma\qsixdel
          =\tfrac{\qsixeps}{2}+\lfloor\qsixeps n/42\rfloor\cdot\tfrac{21}{n}
          \le\tfrac{\qsixeps}{2}+\tfrac{\qsixeps n}{42}\cdot\tfrac{21}{n}
          =\tfrac{\qsixeps}{2}+\tfrac{\qsixeps}{2}=\qsixeps .
\]
Because $|S|=\sigma+1$, Claim~\ref{cl:rank-barrier} applies: $\qsixtL_S$ has at most
$\sigma$ nonzero eigenvalues, and the finiteness of
$\qsixBarrier{u_\sigma}{\sigma}(S)$ forces each of them to lie strictly below
$u_\sigma$. Hence $\qsixopnorm{\qsixtL_S}_{\mathrm{op}}\le u_\sigma\le\qsixeps$, and by
the reduction \eqref{eq:opnorm-form} of Lemma~\ref{lem:opnorm-form} this is
equivalent to $\qsixeps L\succeq L_S$, i.e.\ $S$ is $\qsixeps$-light.

\medskip
Properties (i) and (ii) were both established for the \emph{one} set $S=S_\sigma$
produced by the construction; the lightness is the terminal value of the barrier
invariant and the size is the length of the chain, and these refer to the same
$S$. This proves the theorem with $|S|\ge\lfloor\qsixeps n/42\rfloor$, i.e.\
$c=1/42$.\qed

\section*{Auxiliary lemmas}

\begin{qsixlemma}[Barrier update]\label{lem:bu}
Let $A,B\succeq 0$ with $\qsixrank(B)\le\sigma$ and $\qsixdel>0$. Set $M:=(u+\qsixdel)I-A$. If
$\qsixBarrier{u}{\sigma}(A)<\infty$ and
\[
  \frac{\qsixTrs(M^{-2}B)}{\qsixBarrier{u}{\sigma}(A)-\qsixBarrier{u+\qsixdel}{\sigma}(A)}
   + \qsixTrs(M^{-1}B) \;<\; 1,
\]
then $\qsixBarrier{u+\qsixdel}{\sigma}(A+B)\le\qsixBarrier{u}{\sigma}(A)$.

\smallskip
\noindent\textup{(\emph{Convention.}} Throughout this note---already in the
statement above, in \eqref{eq:U}, in \eqref{eq:total}, and in the proof of
Lemma~\textup{\ref{lem:greedy}}---every top-$\sigma$ trace of a mixed product
$\qsixTrs(A'N)$ in which $A'$ is symmetric and $N\succeq0$ is read in the
\emph{symmetrized sense}
\[
  \qsixTrs(A'N) \;:=\; \qsixTrs\!\bigl(N^{1/2}A'N^{1/2}\bigr),
\]
the right-hand side being the top-$\sigma$ eigenvalue sum of the symmetric matrix
$N^{1/2}A'N^{1/2}$.  This is needed because $A'N$ is generally not symmetric, while
$N^{1/2}A'N^{1/2}$ always is; it is consistent with the ordinary trace, since
cyclicity gives $\qsixTr(A'N)=\qsixTr(N^{1/2}A'N^{1/2})$.  In particular, for the
mixed products appearing here, with $B=CC^{\!\top}$ \textup{(}any factorization,
e.g.\ $C=B^{1/2}$\textup{)},
\begin{equation}\label{eq:bu-shared-spectrum}
  \qsixTrs(M^{p}B) \;=\; \qsixTrs\!\bigl(B^{1/2}M^{p}B^{1/2}\bigr)
                 \;=\; \qsixTrs\!\bigl(C^{\!\top}M^{p}C\bigr),
  \qquad p\in\{-1,-2\},
\end{equation}
because $B^{1/2}M^{p}B^{1/2}$ and $C^{\!\top}M^{p}C$ have the same nonzero
spectrum.  Indeed, set $H:=M^{p/2}$ \textup{(}symmetric and $\succ0$ since
$M\succ0$\textup{)}, $J:=B^{1/2}H$, and $K:=HC$.  Using the elementary
shared-nonzero-spectrum identity ``$XY$ and $YX$ have equal nonzero eigenvalues
with multiplicities'' for $X=B^{1/2}$, $Y=HHB^{1/2}=M^{p}B^{1/2}$ in the first
step and for $X=K^{\!\top}$, $Y=K$ in the last,
\[
  B^{1/2}M^{p}B^{1/2}=JJ^{\!\top},\quad
  J^{\!\top}J=HBH=M^{p/2}BM^{p/2}=HCC^{\!\top}H=KK^{\!\top},\quad
  K^{\!\top}K=C^{\!\top}M^{p}C,
\]
so $B^{1/2}M^{p}B^{1/2}=JJ^{\!\top}$ and $C^{\!\top}M^{p}C=K^{\!\top}K$ share the
nonzero spectrum of the common matrix $J^{\!\top}J=KK^{\!\top}$.  Their top-$\sigma$
eigenvalue sums therefore coincide \textup{(}both equal the sum of those nonzero
eigenvalues, of which there are at most $\qsixrank(B)\le\sigma$, with the lists
padded by zeros\textup{)}.
This identity is what makes the symbol $\qsixTrs(M^{p}B)$ in the hypothesis above
denote the same number as the quantity bounded by the Sherman--Morrison--Woodbury
computation in the proof.\textup{)}
\end{qsixlemma}

\begin{proof}
Throughout, $\qsixTrs$ applied to a product of a symmetric matrix and a PSD
matrix is interpreted with the symmetrized convention stated with the lemma:
for symmetric $A'$ and PSD $N$, $\qsixTrs(A'N):=\qsixTrs(N^{1/2}A'N^{1/2})$, and
the shared-nonzero-spectrum identity~\eqref{eq:bu-shared-spectrum} is used to pass
between $\qsixTrs(M^{p}B)$, $\qsixTrs(B^{1/2}M^{p}B^{1/2})$ and
$\qsixTrs(C^{\!\top}M^{p}C)$.
We also repeatedly use the elementary fact
\begin{equation}\label{eq:rank-collapse}
  \text{if a symmetric matrix }Z\text{ has }\qsixrank(Z)\le\sigma,\text{ then }
  \qsixTrs(Z)=\qsixTr(Z),
\end{equation}
because a symmetric matrix of rank $\le\sigma$ has at most $\sigma$ nonzero
eigenvalues, so the top-$\sigma$ sum already includes every eigenvalue that can
be nonzero; the omitted eigenvalues are all zero. (When $Z\succeq0$ this is the
inequality $\qsixTrs(Z)\le\qsixTr(Z)$ with equality.)

\medskip\noindent\emph{Step 0: invertibility and the basic spectral facts.}\quad
$\qsixBarrier{u}{\sigma}(A)<\infty$ gives $u>\lambda_{\max}(A)$, hence
$M=(u+\qsixdel)I-A\succ0$ and $M^{-1},M^{-2}\succ0$. Write $B=CC^{\!\top}$ with $C$
having $\qsixrank(C)=\qsixrank(B)\le\sigma$ columns (e.g.\ a thin square-root
factorization), so that every matrix of the form $C^{\!\top}(\cdot)C$ has rank
$\le\sigma$. Set
\[
  X:=C^{\!\top}M^{-1}C\succeq0,\qquad Y:=C^{\!\top}M^{-2}C\succeq0,
\]
both of size at most $\sigma\times\sigma$, hence $\qsixrank(X),\qsixrank(Y)\le\sigma$.

\medskip\noindent\emph{Step 1: the hypothesis controls $X$.}\quad
By the symmetrized convention and $\qsixrank(M^{-1/2}BM^{-1/2})\le\qsixrank(B)\le\sigma$
together with \eqref{eq:rank-collapse},
\[
  \qsixTrs(M^{-1}B)=\qsixTrs(B^{1/2}M^{-1}B^{1/2})=\qsixTr(B^{1/2}M^{-1}B^{1/2})
   =\qsixTr(M^{-1}B)=\qsixTr(C^{\!\top}M^{-1}C)=\qsixTr(X),
\]
using cyclicity of the ordinary trace and $B=CC^{\!\top}$. The hypothesis implies
in particular $\qsixTrs(M^{-1}B)<1$, i.e.\ $\qsixTr(X)<1$. Because $X\succeq0$,
\begin{equation}\label{eq:opnorm-trace}
  \qsixopnorm{X}_{\mathrm{op}}=\lambda_{\max}(X)\le\sum_i\lambda_i(X)=\qsixTr(X)<1,
\end{equation}
the inequality being valid since all eigenvalues of a PSD matrix are
non-negative. Hence $I-X\succ0$ is invertible and $(I-X)^{-1}\succeq0$.

\medskip\noindent\emph{Step 2: Sherman--Morrison--Woodbury and a rank-$\sigma$
correction.}\quad
Since $\lambda_{\max}(X)<1$, the matrix $M-CC^{\!\top}=(u+\qsixdel)I-(A+B)$ is
invertible, and the Woodbury identity gives
\begin{equation}\label{eq:woodbury}
  \bigl((u+\qsixdel)I-(A+B)\bigr)^{-1}
   =(M-CC^{\!\top})^{-1}
   = M^{-1} + \underbrace{M^{-1}C(I-X)^{-1}C^{\!\top}M^{-1}}_{=:N}.
\end{equation}
Here $N=M^{-1}C(I-X)^{-1}C^{\!\top}M^{-1}$ is symmetric and PSD (it is
$R(I-X)^{-1}R^{\!\top}$ with $R:=M^{-1}C$ and $(I-X)^{-1}\succeq0$), and
$\qsixrank(N)\le\qsixrank(C)\le\sigma$.

The right-hand side of \eqref{eq:woodbury} is PSD (a sum of PSD matrices),
so $(u+\qsixdel)I-(A+B)\succ0$, i.e.\ $\lambda_{\max}(A+B)<u+\qsixdel$ and
$\qsixBarrier{u+\qsixdel}{\sigma}(A+B)<\infty$. Moreover
$\qsixBarrier{u+\qsixdel}{\sigma}(A+B)=\qsixTrs\bigl((u+\qsixdel)I-(A+B)\bigr)^{-1}
=\qsixTrs(M^{-1}+N)$.

\medskip\noindent\emph{Step 3: split off the correction by Ky Fan
subadditivity.}\quad
$M^{-1}$ and $N$ are both symmetric, so Ky Fan subadditivity of the top-$\sigma$
eigenvalue sum \cite[Thm.~4.3.47a]{HJ12} gives
\begin{equation}\label{eq:kyfan-split}
  \qsixBarrier{u+\qsixdel}{\sigma}(A+B)
   =\qsixTrs(M^{-1}+N)
   \le \qsixTrs(M^{-1})+\qsixTrs(N).
\end{equation}

\medskip\noindent\emph{Step 4: evaluate $\qsixTrs(N)$ via the rank bound and
cyclicity (Step~A of the issue).}\quad
Since $N$ is symmetric with $\qsixrank(N)\le\sigma$, \eqref{eq:rank-collapse}
gives $\qsixTrs(N)=\qsixTr(N)$. Now ordinary cyclicity of the trace applies to
the (square) ordinary trace:
\[
  \qsixTr(N)=\qsixTr\!\bigl(M^{-1}C(I-X)^{-1}C^{\!\top}M^{-1}\bigr)
   =\qsixTr\!\bigl((I-X)^{-1}\,C^{\!\top}M^{-2}C\bigr)
   =\qsixTr\!\bigl((I-X)^{-1}Y\bigr),
\]
where we cycled $C^{\!\top}M^{-1}\cdot M^{-1}C=C^{\!\top}M^{-2}C=Y$ to the front
(note $(I-X)^{-1}$ and $Y$ are square of size $\le\sigma$). The product
$(I-X)^{-1}Y$ has rank $\le\qsixrank(Y)\le\sigma$, and although it need not be
symmetric, its ordinary trace equals the ordinary trace of the \emph{symmetric}
matrix $(I-X)^{-1/2}Y(I-X)^{-1/2}$ by cyclicity:
\[
  \qsixTr\!\bigl((I-X)^{-1}Y\bigr)
  =\qsixTr\!\bigl((I-X)^{-1/2}Y(I-X)^{-1/2}\bigr).
\]
The matrix $W:=(I-X)^{-1/2}Y(I-X)^{-1/2}$ is symmetric PSD with
$\qsixrank(W)\le\qsixrank(Y)\le\sigma$, so by \eqref{eq:rank-collapse} again,
\begin{equation}\label{eq:N-as-W}
  \qsixTrs(N)=\qsixTr(N)=\qsixTr(W)=\qsixTrs(W)
   =\qsixTrs\!\bigl((I-X)^{-1/2}Y(I-X)^{-1/2}\bigr).
\end{equation}
By the symmetrized convention this last quantity is exactly
$\qsixTrs\!\bigl((I-X)^{-1}Y\bigr)$, the object named in the issue.

\medskip\noindent\emph{Step 5: the PSD bound $(I-X)^{-1}\preceq(1-\qsixTr X)^{-1}I$
(Step~B of the issue).}\quad
$X\succeq0$ has eigenvalues $0\le\mu_k\le\lambda_{\max}(X)<1$; the corresponding
eigenvalues of $(I-X)^{-1}$ are $(1-\mu_k)^{-1}$. From \eqref{eq:opnorm-trace},
$\lambda_{\max}(X)\le\qsixTr(X)<1$, hence
$1-\qsixTr(X)\le 1-\lambda_{\max}(X)\le 1-\mu_k$, and since all three quantities
are positive, inverting reverses the order:
\[
  \frac{1}{1-\mu_k}\le\frac{1}{1-\lambda_{\max}(X)}\le\frac{1}{1-\qsixTr(X)}
  \qquad\text{for every }k.
\]
Thus every eigenvalue of $(I-X)^{-1}$ is at most $(1-\qsixTr X)^{-1}$, i.e.
\begin{equation}\label{eq:invX-bound}
  (I-X)^{-1}\preceq \frac{1}{1-\qsixTr(X)}\,I .
\end{equation}
(This is the step the issue flags as counter-intuitive: replacing the small
operator norm $\lambda_{\max}(X)$ by the larger quantity $\qsixTr(X)$ makes the
\emph{denominator smaller} and hence the bound \emph{larger}, so the
inequality \eqref{eq:invX-bound} is valid.)

\medskip\noindent\emph{Step 6: conclude via PSD-monotonicity of $\qsixTrs$
(Step~C of the issue).}\quad
Conjugating the PSD inequality \eqref{eq:invX-bound} by $Y^{1/2}$ preserves the
order:
\[
  Y^{1/2}(I-X)^{-1}Y^{1/2}\;\preceq\;\frac{1}{1-\qsixTr(X)}\,Y^{1/2}IY^{1/2}
   =\frac{1}{1-\qsixTr(X)}\,Y .
\]
The top-$\sigma$ eigenvalue sum $\qsixTrs$ is monotone with respect to the PSD
order on symmetric matrices (Ky Fan: $P\preceq Q\Rightarrow\lambda_i(P)\le
\lambda_i(Q)$ for all $i$ by Weyl monotonicity, hence the partial sums obey the
same inequality), so, using also $Y^{1/2}(I-X)^{-1}Y^{1/2}\succeq0$ and the
symmetrized convention,
\[
  \qsixTrs\!\bigl((I-X)^{-1}Y\bigr)
   =\qsixTrs\!\bigl(Y^{1/2}(I-X)^{-1}Y^{1/2}\bigr)
   \le \frac{1}{1-\qsixTr(X)}\,\qsixTrs(Y).
\]
Finally $\qsixrank(Y)\le\sigma$ gives, by \eqref{eq:rank-collapse} and
cyclicity, $\qsixTrs(Y)=\qsixTr(Y)=\qsixTr(C^{\!\top}M^{-2}C)=\qsixTr(M^{-2}B)
=\qsixTr(B^{1/2}M^{-2}B^{1/2})=\qsixTrs(M^{-2}B)$ (the last equality again using
$\qsixrank(B^{1/2}M^{-2}B^{1/2})\le\sigma$ and the symmetrized convention).
Likewise $\qsixTr(X)=\qsixTrs(M^{-1}B)$ as in Step~1. Substituting these and the
result of Step~4,
\begin{equation}\label{eq:N-final}
  \qsixTrs(N)=\qsixTrs\!\bigl((I-X)^{-1}Y\bigr)
   \le \frac{\qsixTrs(M^{-2}B)}{1-\qsixTrs(M^{-1}B)} .
\end{equation}

\medskip\noindent\emph{Step 7: assemble.}\quad
Combining \eqref{eq:kyfan-split} and \eqref{eq:N-final},
\[
  \qsixBarrier{u+\qsixdel}{\sigma}(A+B)
   \le \qsixTrs(M^{-1})+\frac{\qsixTrs(M^{-2}B)}{1-\qsixTrs(M^{-1}B)} .
\]
Now $\qsixTrs(M^{-1})=\qsixBarrier{u+\qsixdel}{\sigma}(A)$ (indeed
$M^{-1}=((u+\qsixdel)I-A)^{-1}$, and $\qsixBarrier{u+\qsixdel}{\sigma}(A)
=\qsixTrs(((u+\qsixdel)I-A)^{-1})$ by definition), so writing
$D:=\qsixBarrier{u}{\sigma}(A)-\qsixBarrier{u+\qsixdel}{\sigma}(A)\ge0$ we have
$\qsixTrs(M^{-1})=\qsixBarrier{u}{\sigma}(A)-D$. The hypothesis reads
\[
  \frac{\qsixTrs(M^{-2}B)}{D}+\qsixTrs(M^{-1}B)<1
  \quad\Longleftrightarrow\quad
  \frac{\qsixTrs(M^{-2}B)}{1-\qsixTrs(M^{-1}B)}<D ,
\]
where the rearrangement is legitimate because $D>0$ (else the first fraction is
$+\infty$ and the hypothesis fails) and $1-\qsixTrs(M^{-1}B)=1-\qsixTr(X)>0$ by
Step~1. Therefore
\[
  \qsixBarrier{u+\qsixdel}{\sigma}(A+B)
   \le \bigl(\qsixBarrier{u}{\sigma}(A)-D\bigr)
       +\frac{\qsixTrs(M^{-2}B)}{1-\qsixTrs(M^{-1}B)}
   < \bigl(\qsixBarrier{u}{\sigma}(A)-D\bigr)+D
   = \qsixBarrier{u}{\sigma}(A),
\]
which is the asserted bound (in fact with strict inequality).
\end{proof}


\begin{qsixclaim}\label{cl:phi-diff}
$\qsixBarrier{u}{\sigma}(S)-\qsixBarrier{u+\qsixdel}{\sigma}(S)\ge\qsixdel\,\qsixTrs(M^{-2})$.
\end{qsixclaim}
\begin{proof}
For each top-$\sigma$ eigenvalue $\lambda_i$ of $\qsixtL_S$,
\[
  \tfrac{1}{u-\lambda_i}-\tfrac{1}{u+\qsixdel-\lambda_i}
   = \tfrac{\qsixdel}{(u-\lambda_i)(u+\qsixdel-\lambda_i)}
   \ge \tfrac{\qsixdel}{(u+\qsixdel-\lambda_i)^2};
\]
sum over $i$.
\end{proof}

\section*{Remarks}

\begin{enumerate}\setlength\itemsep{4pt}
\item \textbf{Constant.} The constant $1/42$ comes from the parameter choices
$\qsixdel=21/n$, $\phi=n/21$, which balance the $5/\qsixdel$ and $5\phi$ contributions
in \eqref{eq:total}. It is not believed optimal; improving it is open.
\item \textbf{Key innovation.} The standard upper-barrier function tracks all
eigenvalues; here we track only the top $\sigma$. This adaptation, together
with the Ky Fan inequalities for $\qsixTrs$, is what makes the argument go
through.
\item \textbf{Connection to spectral sparsification.} An $\qsixeps$-light $S$ with
$|S|\ge c\qsixeps n$ says the induced subgraph on $S$ contributes at most $\qsixeps L$
spectrally, giving a perspective on ``light'' as ``spectrally independent''.
\item \textbf{Vertex vs.\ edge sampling.} The edge-partitioning analogue
admits an optimal constant of $1/2$ via the Kadison--Singer resolution. The
vertex-light variant is more delicate because edges sharing a vertex are
correlated; the deterministic barrier method bypasses concentration entirely.
\end{enumerate}

\begin{thebibliography}{9}
\bibitem{HJ12} R.~A.~Horn and C.~R.~Johnson.
  \emph{Matrix Analysis}, 2nd ed.\ Cambridge University Press, 2012.
\bibitem{Bhatia97} R.~Bhatia.
  \emph{Matrix Analysis}, Graduate Texts in Mathematics 169.\ Springer, 1997.
\end{thebibliography}

\endgroup